\color{unidarkgreen}{\subsection{Контроль синхронизма и напряжений}}
\color{black}

\begin{enumerate}[label=\arabic{section}.\arabic{subsection}.\arabic{subsubsection}.\arabic*,  align=left, labelsep=4pt, leftmargin=0pt, itemindent=57pt]

\item
Функциональный блок контроля синхронизма и напряжений (КСН) предназначен для формирования разрешающих сигналов, соответствующих условиям включения, для АПВ, оперативного включения. Функциональный блок КСН В1(ЛВ)СН аналогичен КСН В1(ЛВ)ВН.

\item
Основные режимы включения и контролируемые для них условия приведены в таблице \ref{desc:tbl1}. Режимы включения с контролем синхронизма (КС) и с улавливанием синхронизма (УС) вводятся соответствующими программными ключами \textbf{«КС В1(ЛВ) ВН»} и \textbf{«УС В1(ЛВ) ВН»}.

\item
Если разность частот напряжений шин и трансформатора не превышает уставку \textbf{«df макс КС»}, то разрешается включение с контролем синхронизма. Для режима включения с КС также контролируется:
\begin{itemize}
\item наличие напряжения на обоих энергообъектах (шины и трансформатор);
\item разность напряжений шин и трансформатора, не превышающая по модулю значения заданной уставки;
\item разность фаз напряжений шин и трансформатора, не превышающая по модулю значения заданной уставки.
\end{itemize}

\item
При выводе программного ключа \textbf{«КС В1(ЛВ) ВН»} включение разрешается без контроля частоты, а только по уровню напряжений и разности их фаз.

\item
Если разность частот превышает уставку \textbf{«df макс КС»}, но не выходит за пределы уставки \textbf{«df пред. доп.»}, то разрешается включение с улавливанием синхронизма. Сигнал включения с улавливанием синхронизма формируется с опережением на время включения выключателя. Режим включения с УС можно вывести программным ключом \textbf{«УС В1(ЛВ) ВН»}.

\item
Наличие или отсутствие напряжения на трансформатор контролируется измерительными органами максимального и минимального напряжения, пороги срабатывания которых задаются уставками \textbf{«Uнал Т В1(ЛВ) ВН»} и \textbf{«Uотс Т В1(ЛВ) ВН»} соответственно. Пороги срабатывания измерительных органов для контроля напряжения на энергообъекте Ш (шины) задаются уставками \textbf{«Uнал Ш В1(ЛВ) ВН»} и \textbf{«Uотс Ш В1(ЛВ) ВН»}.

\small
\begin{longtable}{|>{\centering\arraybackslash}m{5.3cm}|>{\centering\arraybackslash}m{11.6cm}|}
% Первая страница: полный заголовок
\caption{Условия включения\hfill\vspace{-0.5\baselineskip}}\label{desc:tbl1}\\
\hline
\rowcolor{gray!30}
Режим включения & Условия включения \\
\hline
\endfirsthead

% Промежуточные страницы: «Продолжение таблицы»
\caption*{\hspace{3pt}\emph{Продолжение таблицы \ref{desc:tbl1}}\hfill\ }\vspace{-0.5\baselineskip} \\
\hline
\rowcolor{gray!30}
Режим включения & Условия включения \\
\hline
\endhead

ОНАТ+ННШ &   
			\begin{itemize}
			\vspace{3pt}
			\item $ \textrm{U}_\textrm{T} $ меньше уставки \textbf{«Uотс АТ В1(ЛВ) ВН»};
			\item $ \textrm{U}_\textrm{Ш} $ больше уставки \textbf{«Uнал Ш В1(ЛВ) ВН»}.
			\vspace{3pt}
			\end{itemize} \\
\hline
ННШ+ННАТ & 
			\begin{itemize}
			\vspace{3pt}
			\item $ \textrm{U}_\textrm{T} $ больше уставки \textbf{«Uнал АТ В1(ЛВ) ВН»};
			\item $ \textrm{U}_\textrm{Ш} $ меньше уставки \textbf{«Uотс Ш В1(ЛВ) ВН»}.
			\vspace{3pt}
			\end{itemize} \\
\hline
КС(УС)	&
			\vspace{3pt}
			
			\begin{itemize}
			\item[] \noindent 1) Включение контролем синхронизма:
			\item $ \textrm{U}_\textrm{T} $ больше уставки \textbf{«Uнал АТ В1(ЛВ) ВН»};
			\item $ \textrm{U}_\textrm{Ш} $ больше уставки \textbf{«Uнал Ш В1(ЛВ) ВН»};
			\item разность напряжений трансформатора и шин не превышает по модулю значения заданной уставки \textbf{«dU макс КС(УС)»};
			\item разность фаз напряжений трансформатора и шин не превышает по модулю значения заданной уставки \textbf{«dФ макс КС»};
			\item разность частот напряжений трансформатора и шин не превышает значения заданной уставки \textbf{«df макс КС»};
			\end{itemize}
			
			\begin{itemize}
			\item[] \noindent  2) Включение с улавливанием синхронизма:
			\item $ \textrm{U}_\textrm{T} $ больше уставки \textbf{«Uнал АТ В1(ЛВ) ВН»};
			\item $ \textrm{U}_\textrm{Ш} $ больше уставки \textbf{«Uнал Ш В1(ЛВ) ВН»};
			\item разность напряжений трансформатор и шин не превышает по модулю значения заданной уставки \textbf{«dU макс КС(УС)»};
			\item разность частот напряжений трансформатора и шин превышает значение уставки \textbf{«df макс КС»};
			\item разность частот напряжений трансформатора и шин не превышает значения заданной уставки \textbf{«df пред. доп.»};
			\item за время \textbf{«Т опер. вкл. В1(ЛВ) ВН»} с момента подачи сигнала на включение выключателя расхождение углов между векторами напряжений на трансформаторе и шине превысит 10º.
			\vspace{3pt}
			\end{itemize} \\
\hline

\end{longtable}
\normalsize	

\item
Сигнал разрешения АПВ автотрансформатор \textit{«Разр. АПВат В1(ЛВ) ВН»} формируется при выполнении условий включения в режимах ОНАТ+ННШ, КС или УС.

\item
Сигнал разрешения АПВ шин \textit{«Разр. АПВш В1(ЛВ) ВН»} формируется при выполнении условий включения в режимах ОНШ+ННАТ, КС или УС. 

\item
Сигнал разрешения оперативного включения \textit{«Разр. опер. вкл. В1(ЛВ) ВН»} может формироваться от внутренней логики КСН или от внешнего сигнала в зависимости от положения программного переключателя \textbf{«Контр. опер. вкл. В1(ЛВ) ВН»}.

\item
В положении \textit{«Внутр.»} оперативное включение разрешается при выполнении условий включения в режимах ОНАТ+ННШ (при опробовании автотрансформатора), ОНШ+ННАТ (при опробовании шин), КС или УС (при замыкании в транзит). Также предусмотрена возможность оперативного включения с контролем отсутствия напряжения на шинах или на автотрансформатор при введенном программном ключе \textbf{«Опер. вкл. В1(ЛВ) ВН с КОН»}. При отсутствии внешнего сигнала \textit{«Опер. вкл. с КС(УС) В1(ЛВ) ВН»} оперативное включение может выполняться без контроля напряжений и условий синхронизма.

\item
В положении \textit{«Внеш.»} оперативное включение разрешается при наличии внешнего сигнала \textit{«Внеш. сигнал КС(УС) В1(ЛВ) ВН»}.

\item
Сигнал разрешения полуавтоматического включения \textit{«Разр. ПАВ В1(ЛВ) ВН»} формируется при выполнении условий включения в режимах КС или УС.

\item
Алгоритм работы КСН ВН приведен на рисунке \ref{pic:KSN}, алгоритм работы КСН СН выполнен аналогично. Уставки КСН ВН приведены в таблице \ref{app_set:ksnvn}. Уставки КСН СН приведены в таблице \ref{app_set:ksnsn}.

\begin{figure}[H]
\centering
\includegraphics[width=0.8\textwidth,height=0.8\textheight,keepaspectratio]{\fbpath/761.3351 КСН/img/KSN.pdf}
\captionof{figure}{Функциональная схема КСН В1(ЛВ) ВН}\label{pic:KSN}
\end{figure}

\end{enumerate}